\فصل{مقدمه}
\label{ch:introduction}

% Keep the existing bibliography and citation numbering after replacing the template prose.
\nocite{knuth1984texbook,lamport1985LaTeX}

\قسمت{تعریف مسئله و انگیزه}

حرکت ذرات در جریان لزج کانالی به میدان سرعت سیال وابسته است.
مکان و سرعت اولیهٔ ذره نیز بر مسیر آن اثر می‌گذارند.
زمان رهاسازی بر موقعیت ذرات در زمان هدف اثر می‌گذارد.
در مسئلهٔ مغناطیسی، تحریک منابع نیز مسیر ذره را تغییر می‌دهد.

در مسئلهٔ مستقیم، میدان یا مسیر از ورودی معلوم پیش‌بینی می‌شود.
برای نمونه، پروفایل ورودی معلوم می‌تواند مبنای پیش‌بینی میدان جریان باشد.
در بازسازی معکوس جریان، مسیر مشاهده‌شده برای تخمین پروفایل ورودی به‌کار می‌رود.
در طراحی معکوس مغناطیسی، تحریک مجاز برای بازتولید مسیر هدف تعیین می‌شود.

طراحی رهاسازی و تحریک، معمولاً به ارزیابی چندین نامزد نیاز دارد.
مدل یادگرفته‌شده امکان ارزیابی مکرر نامزدها را فراهم می‌کند.
در بازسازی جریان، مسیر ذره فقط بخشی از کانال را پوشش می‌دهد.
پوشش محدود مسیرها، اطلاعات موجود برای بازسازی را کاهش می‌دهد.
در برخی مشاهدات، تنها هندسهٔ مسیر در دسترس است.
در این حالت، زمان‌بندی حرکت نیز در برازش تخمین زده می‌شود.

\قسمت{اهداف و سهم‌های پژوهش}

\begin{enumerate}
\item
شبکهٔ عصبی فیزیک‌آگاه پارامتریک، مسیر تک‌ذره را از شرایط رهاسازی پیش‌بینی می‌کند.
این پیش‌بینی برای طراحی مکان، سرعت و زمان رهاسازی ذرات به‌کار می‌رود.
بازپخش مستقل، آرایش هدف و فاصلهٔ ذرات هم‌زمان حاضر را بررسی می‌کند.

\item
عملگر عمیق (\lr{DeepONet})، میدان جریان را از پروفایل ورودی پیش‌بینی می‌کند.
عملگر مغناطیسی نیز مسیر ذره را از تحریک معلوم دو منبع می‌سازد.

\item
پروفایل ورودی از مسیرهای زمان‌دار یا هندسهٔ مسیر بازسازی می‌شود.
اثر تعداد ذرات و محل رهاسازی بر خطای پروفایل بررسی می‌شود.
دقت بازپخش مسیر و دقت پروفایل یا میدان جداگانه سنجیده می‌شوند.

\item
برای بازتولید مسیر هدف، یک زوج تحریک مغناطیسی مجاز طراحی می‌شود.
حل‌گر مستقل، مسیر حاصل از تحریک طراحی‌شده را محاسبه می‌کند.
\end{enumerate}

\قسمت{حدود پژوهش}

مدل‌ها تحت فرض‌های فیزیکی و دامنه‌های فصل~\ref{ch:theory} بررسی می‌شوند.
میدان جریان با پاسخ اجزای محدود مقایسه می‌شود.
در آزمون انتقال مغناطیسی منتشرشده، پاسخ تحلیلی مرجع مقایسه است.
طرح‌های رهاسازی و تحریک با بازپخش عددی مستقل ارزیابی می‌شوند.

پیشینهٔ شبکه‌های فیزیک‌آگاه و یادگیری عملگر، زمینهٔ انتخاب این روش‌ها را فراهم می‌کند.
فصل~\ref{ch:literature-review} مطالعات مربوط به هر مسئله را مرور می‌کند.

\قسمت{ساختار پایان‌نامه}

فصل~\ref{ch:introduction}، «مقدمه»، مسئله و اهداف پژوهش را معرفی می‌کند.
فصل~\ref{ch:literature-review}، «مروری بر ادبیات»، پیشینهٔ مرتبط را بررسی می‌کند.
فصل~\ref{ch:theory}، «تئوری مسئله»، مدل‌ها و روش‌های حل را صورت‌بندی می‌کند.
فصل~\ref{chap:results}، «نتایج»، آزمون‌ها و ارزیابی‌ها را گزارش می‌کند.
فصل~\ref{ch:conclusions}، «نتیجه‌گیری»، یافته‌ها، محدودیت‌ها و مسیرهای آتی را جمع‌بندی می‌کند.
